\begin{algorithm}
	\caption{Effective Degree}
	\label{alg:effective-deg}
	\begin{algorithmic}
		\item on Initialization: for each offline vertex $u_i$ let $d_{u_i}=\deg(u_i)$
		\item on Arrival of vertex $v$: $N$ are unmatched offline neighbors of $v$. If $N\neq\emptyset$, match $v$ to $\underset{u\in N}{\mathrm{argmin}}\ d_u$, $\underset{u\in N}{\forall}d_u-=1$
	\end{algorithmic}
\end{algorithm}
Despite this algorithm seeming promising, it is only marginally more effective than a naive deterministic approach. We will show that in this model no deterministic algorithm has asymptocic competitiveness greater then $0.5$
\begin{theorem}
	\label{thm:deg-det-comp}
	In the Final Degree model no deterministic algorithm has a competitiveness higher than $0.5$
\end{theorem}
\begin{proof}
	Consider the graph $G_{\deg}$ which has $n=2k+1$ offline vertices with final degree $k+1$ and is formed as follows:\\
	$k+1$ vertices arrive, all connected to all unmatched offline vertices.\\
	Notice that at this point $k+1$ vertices are matched and the remaining $k$ have degree $k+1$ in the current graph, meaning that they cannot be matched. All that remains is to fulfill the degree promise. This is easy as the degrees of the matched vertices are $1,2,\dots,k,k+1$ meaning that all that is required is for $k$ vertices connected to all offline vertices with unfulfilled degree promise to arrive.\\
	The size of the matching found is $k+1$ while the graph has a perfect matching of size $n=2k+1$ which is achived when the $k$ vertices matched by the graph are instead matched to the vertices which fulfilled their degree promise while the remaining $k+1$ vertices are matched among themselves since they form a biclique $K_{k+1,k+1}$
\end{proof}
\begin{figure}[H]
	\label{fig:deg-det-comp}
	\centering
	\scalebox{.8}{\begin{tikzpicture}
	\begin{pgfonlayer}{nodelayer}
		\node [style=default] (0) at (-3, 13) {};
		\node [style=default] (1) at (-3, 12) {};
		\node [style=default] (2) at (-3, 11) {};
		\node [style=default] (3) at (-3, 10) {};
		\node [style=default] (4) at (-3, 9) {};
		\node [style=default] (5) at (-3, 8) {};
		\node [style=default] (6) at (-3, 7) {};
		\node [style=deleted] (7) at (2, 13) {};
		\node [style=deleted] (8) at (2, 12) {};
		\node [style=deleted] (9) at (2, 11) {};
		\node [style=deleted] (10) at (2, 10) {};
		\node [style=deleted] (11) at (2, 9) {};
		\node [style=deleted] (12) at (2, 8) {};
		\node [style=deleted] (13) at (2, 7) {};
		\node [style=none] (14) at (-4, 13) {4};
		\node [style=none] (15) at (-4, 12) {4};
		\node [style=none] (16) at (-4, 11) {4};
		\node [style=none] (17) at (-4, 10) {4};
		\node [style=none] (18) at (-4, 9) {4};
		\node [style=none] (19) at (-4, 8) {4};
		\node [style=none] (20) at (-4, 7) {4};
		\node [style=none] (21) at (2.75, 10) {1};
		\node [style=none] (22) at (2.75, 9) {2};
		\node [style=none] (23) at (2.75, 8) {3};
		\node [style=none] (24) at (2.75, 7) {4};
		\node [style=none] (25) at (2.75, 11) {5};
		\node [style=none] (26) at (2.75, 12) {6};
		\node [style=none] (27) at (2.75, 13) {7};
		\node [style=none] (28) at (-3.25, 6) {final degree};
		\node [style=none] (29) at (2, 6) {arrival order};
	\end{pgfonlayer}
	\begin{pgfonlayer}{edgelayer}
		\draw (6) to (13);
		\draw (5) to (12);
		\draw (4) to (11);
		\draw (3) to (10);
		\draw (2) to (9);
		\draw (1) to (8);
		\draw (0) to (7);
		\draw (0) to (8);
		\draw (0) to (9);
		\draw (1) to (9);
		\draw (1) to (10);
		\draw (2) to (10);
		\draw (2) to (11);
		\draw (3) to (11);
		\draw (3) to (12);
		\draw (4) to (12);
		\draw (4) to (13);
		\draw (4) to (10);
		\draw (5) to (10);
		\draw (5) to (11);
		\draw (5) to (13);
		\draw (6) to (12);
		\draw (6) to (11);
		\draw (6) to (10);
		\draw (0) to (10);
		\draw (1) to (11);
		\draw (2) to (12);
		\draw (3) to (13);
	\end{pgfonlayer}
\end{tikzpicture}
}
\end{figure}
We will remark here that this model is a relaxation of the Semi-Online model, so competitiveness here isn't lower\\

\subsection{Randomized Augmentation}
As shown above, the information of a graph's final degree isn't helpful for deterministic algorithms. However, there is a good reason to believe that this might not be the case for rndomized algorithms\\
Here, the fact that the competitiveness is only $0.5$ asymptotically and in actuality is $\frac{k+1}{2k+1}$ is key.

\begin{lemma}
\label{deg-diff-degree}
If all vertices have different degrees then the algorithm finds a perfect matching
\end{lemma}

\begin{proof}
	Let $G_n=(V,U,E)$ be a $n\times n$ bipartite graph with with all vertices in $U$ having different degrees. WLOG $G$ has a perfect matching, implying that the degrees of $u\in U$ are $1,\dots,n$. Let $u_k:=u\in U|deg(u)=k$\\
	To force the algorithm to leave $n_k$ behind, $k$ online vertices need to arrive, each connected to a different $u'_i\in\{u_i|i<k\}$. The problem with this is that the size of $U'$ is $k-1$, maiing sich a thing impossible
\end{proof}

A much stronger realization can be gleamed from this reasoning to force a deterministic algorithm to leave a vertex $u$ behind, and for it to be possible for a randomized algorithm to leave that vertex behind, there must be $\deg(u)$ vertices $v$ such that $\deg(u)\geq\deg(v)$, which moreover won't be left behind.\\

This line of reasoning leads directly to the graph $G_{\deg}$ used in the proof of \autoref{thm:deg-det-comp}. However, this also leads to a problem for an adversary. That graph is rather dense, which isn't a situation where randomized graphs perform poorly. Indeed, despite
RGA (\autoref{alg:RGA}) having an expected competitiveness of $\frac{n}{2}+O(\log n)$, it is expected to find a matching of size $n(1-\frac{1}{e})$ on a 
%reference
graph which is actually a subgraph of the degree countergraph edgewise, meaning it would work at least as well here

\begin{algorithm}
	\caption{Randomized Effective Degree}
	\label{alg:effective-deg-rand}
	\begin{algorithmic}
		\item on Initialization: for each offline vertex $u_i$ let $d_{u_i}=\deg(u_i)$
		\item on Arrival of vertex $v$: $N$ are unmatched offline neighbors of $v$. If $N\neq\emptyset$, let $d=\underset{u\in N}{\mathrm{min}}\ d_u$. $v=\mathrm{random}(\{u\in N|\deg(v)=d\})$. Match $v$ to $u$, $\underset{u\in N}{\forall}d_u-=1$
	\end{algorithmic}
\end{algorithm}

\begin{lemma}
	On $G_{\deg}$ from \autoref{fig:deg-det-comp} \autoref{alg:effective-deg-rand} has an expected performance of no less than $\approx 0.75$
\end{lemma}
\begin{proof}
	We can divide the arriving online vertices into two phases: first being when the first $k+1$ vertices that make the upper $k$ offline vertices unmatchable in the deterministic scenario, and the second phase when latter $k$ online vertices arrive and are unmatchable.\\
	In the first phase the arriving vertices are always matchable since they are connected to all offline vertices\\
	In the second phase arriving vertices are connected to incrementally smaller parts of the upper $k$ vertices\\
	The lower $k+1$ offline and online vertices form a biclique, meaning that all matches made between the first $k+1$ online vertices and the lower $k+1$ offline vertices do not make any vertex unmatchable. Since matches are chosen randomly, we expect half of the online agends from the first phase to be matched there\\
	In the second phase $k$ vertices arrive connected to incrementally smaller parts of the upper $k$ offline vertices. Using an argument symmetrical to the one in \autoref{deg-diff-degree} we know that the matching here is perfect\\
	In total - all online vertices in the first phase have been matched and an expected half of the vertices in the second phase has ben matched, resulting in an expected performance of $0.75$
\end{proof}

Note here that we have conducted Monte Carlo experiments  with this exact algorithm and graph for $k$ up to $50$ and in over 1 million simulations the algorithm found a matching of average size $n-2$ suggesting that the actual performance is much higher here

<If it is possible to prove that different degrees always result in the algorithm kaing better decisions, then \autoref{alg:effective-deg-rand} has a competitiveness of $0.75$ since minimizing it's ability to make correct decisions, the final graph must be a series of disconnected $G_{\deg}$ of various sizes>